\documentclass{article}
\usepackage{amsmath, amssymb, amsthm}
\usepackage{hyperref}
\usepackage{geometry}
\geometry{margin=1in}

\title{Erd\H{o}s Problem \#536}
\date{Last updated: 2025-08-31}
\author{Source: erdosproblems.com}

\begin{document}
\maketitle

\noindent\textbf{Status:} OPEN \quad \textbf{No prize}

\noindent\textbf{Tags:} number theory

\medskip
\noindent\textbf{Problem Statement:}

Let $f(N)$ be the largest size of $A\subseteq \{1,\ldots,N\}$ with the property that there are no distinct $a,b,c\in A$ such that\[[a,b]=[b,c]=[a,c],\]where $[a,b]$ denotes the least common multiple.

Estimate $f(N)$ - in particular, is it true that $f(N)=o(N)$?

\bigskip
\noindent\textbf{Remarks:}

The corresponding quantity is $\gg N$ if we ask for four elements with the same pairwise least common multiple, as shown by Erd\H{o}s \cite{Er62} (with a proof given in \cite{Er70}).

This was also asked by Erd\H{o}s at the 1991 problem session of West Coast Number Theory.

Abbott and Gardner \cite{AbGa67} proved\[f(N) \geq (1-o(1))(\log\log N)\frac{N}{\log N}.\]In the comments Weisenberg sketches a proof that, for some function $\omega(N)\to \infty$ as $N\to \infty$,\[(\log\log N)^{\omega(N)}\frac{N}{\log N}\ll f(N) \leq \left(\frac{221}{225}+o(1)\right)N.\]See also [535], [537], and [856]. A related combinatorial problem is asked at [857].

\bigskip
\noindent\textbf{References:}

[AbGa67] Abbott, H. L. and Gardner, B., An extremal problem in number theory. Canad. Math. Bull. (1967), 173--177.
 
 [Er62] Erd\H{o}s, P\'{a}l, Remarks on number theory. IV. Extremal problems in number theory. I. Mat. Lapok (1962), 228-255.
 
 [Er70] Erd\H{o}s, Paul, Some extremal problems in combinatorial number theory. Mathematical Essays Dedicated to A. J. Macintyre (1970), 123-133.

\bigskip
\noindent\small{Source: \url{https://www.erdosproblems.com/536}}
\end{document}
